\section{Logical Relations for \delamlang}\labeledit{sec:dt:logrel}

Previously, we have given the judgments of \delamlang, verified its syntactic
properties and given its reduction and convertibility algorithms. %
Starting this section, we establish the logical relation and prove the (weak)
normalization and convertibility properties of \delamlang. %
Following \citet{abel_decidability_2017}, we proceeds as follows:
\begin{itemize}
\item First we give a set of generic equivalence conditions for a parameterized
  discussion of the logical relation.
\item Then we give the definition of the Kripke logical relations of types and
  terms. %
  The logical relations are parameterized by layers. %
  In this step, we are only concerned about types that are available at all layers,
  i.e. those in MLTT and unrelated to meta-programming. 
\item Then we give the definition  of the Kripke logical relations of local contexts
  and local substitutions. %
\item Then we branch off two orthogonal developments.
  \begin{itemize}
  \item We give the definition of the Kripke logical relations of global contexts and
    global substitutions.
  \item We give the definition of the Kripke logical relations of types and terms,
    again. %
    But in this case, we must also give the definition for types that are related to
    meta-programming, i.e. contextual types. %
  \end{itemize}
  In fact, the definitions given by the two sub-steps above must consider each
  other. %
  Otherwise, we will not able to extend related global substitutions during the proof
  of the fundamental theorems. 
\item Next we give the semantic judgments. %
  The semantic judgments require types, terms, etc. to be stable under all universe,
  global and local substitutions. %
\item Finally, we establish the fundamental theorems for the semantic judgments. %
  Instantiating the generic equivalence gives us the proof of convertibility.
\end{itemize}
Due to layering, following \Cref{sec:cv:logrel}, the generic equivalence, logical
relation and validity judgments are all layered. %
In fact, since computation exists at both layers $\proglevel$ and $\metalevel$, the situation is very
complex. %
\citet{abel_decidability_2017} instantiate their generic equivalence twice to obtain
the decidability of convertibility checking, and we will also be doing the same. %
Due to the complication of layering, our fundamental theorems must talk about all
layers. %
The difficulties of the logical relations lie in that how we can support code running
and recursions on code at the same time and justify them in the semantics. %

\subsection{Generic Equivalence}

Similar to \Cref{sec:cv:logrel}, we first quantify four generic equivalence relations,
which will be instantiated to syntactic equivalence and convertibility later, and
their laws. %
This step provides modularity to logical relation argument: we simply instantiate this
generic equivalence to obtain different versions of the fundamental theorems. %
Due to dependent types, we define generic equivalence over $i \in \{\proglevel, \metalevel\}$:
\begin{itemize}
\item $\lttypgneeq i V{V'} l$ describes a generic type equivalence between two neutral
  types at universe level $l$ at layer $i$.
\item $\lttypgeq i {T}{T'} l$ describes a generic type equivalence between two types at
  universe level $l$ at layer $i$.
\item $\lttrmgneeq i \nu{\nu'} T l$ describes a generic type equivalence between two neutral
  terms of type $T$ at universe level $l$ at layer $i$. 
\item $\lttrmgeq i t{t'} T l$ describes a generic type equivalence between two
  terms of type $T$ at universe level $l$ at layer $i$. 
\end{itemize}

From the four generic equivalence, we induce two equivalence of local contexts and
local substitutions by using the generic equivalence pairwise:
\begin{mathpar}
  \inferrule
  {\judge[L] \Psi}
  {\ltctxgeq i \cdot\cdot}

  \inferrule
  {\judge[L] \Psi \\ g : \Ctx \in \Psi}
  {\ltctxgeq i{g}{g}}

  \inferrule
  {\ltctxgeq i \Gamma\Delta \\ \lttypgeq i T{T'} l \\ \tyequiv[L]l{l'}\Level}
  {\ltctxgeq i{\Gamma, x : T \at l}{\Delta, x : T' \at{l'}}}

  \inferrule
  {\lpjudge i\Gamma \\ \text{$\Gamma$ ends with $\cdot$} \\ |\Gamma| = \metalevel}
  {\ltsubgeq i {\cdot^\metalevel}{\cdot^\metalevel}{\cdot}}

  \inferrule
  {\lpjudge i\Gamma \\ g : \Ctx \in \Psi \\ \text{$\Gamma$ ends with $g$} \\ |\Gamma| = \metalevel}
  {\ltsubgeq i {\cdot_g^\metalevel}{\cdot_g^\metalevel}{\cdot}}

  \inferrule
  {\lpjudge i\Gamma \\ g : \Ctx \in \Psi \\\\ \text{$\Gamma$ ends with $g$} \\ |\Gamma| = \metalevel}
  {\ltsubgeq i {\wk_g^\metalevel}{\wk_g^\metalevel}{g}}

  \inferrule
  {\ltsubgeq i {\delta}{\delta'}{\Delta} \\\\ \lttypwf[\Psi][\Delta]i T l \\
    \lttrmgeq i {t}{t'}{T[\delta]} l}
  {\ltsubgeq i {\delta, t/x}{\delta', t'/x}{\Delta, x : T \at l}}
\end{mathpar}

The generic equivalence and the logical relations are invariant under all
weakenings. %
Therefore, we should make these notions clear here. %
Since there are three different contexts, we have three corresponding kinds of
weakenings. %
In particular, $\theta :: L \To L'$ is the universe weakening. %
Following previous conventions, $\gamma :: L \sep \Psi \To_g \theta$ is a global
weakening, and $\tau :: L \sep \Psi;\Gamma \To_i \Delta$ is a local weakening. %
The subscript $i$ denotes which layer the contexts $\Gamma$ and $\Delta$ live in. %
We can simultaneously weaken all three contexts at the same time. %
We simply apply $\theta$, $\gamma$ and $\tau$ in this order. %
We let $\psi$ represent this triple:
\[
  \psi := (\theta, \gamma, \tau) :: L \sep \Psi; \Gamma \To_i L' \sep \Phi;\Delta
\]
where
\begin{align*}
  \theta &:: L \To L' \\
  \gamma &:: L' \sep \Psi[\theta] \To_g \Phi \\
  \tau &:: L' \sep \theta; \Gamma[\theta][\gamma] \To_i \Delta
\end{align*}
Similarly we let
\[
  \alpha := (\theta, \gamma) :: L \sep \Psi \To L' \sep \Phi
\]
We can apply weakenings like substitutions to universe levels, types, terms, contexts
and substitutions as expected. %
The action is to shift the variables according to the specified weakenings. %
This is a standard action, despite having three separate notions, so we take it for
granted here. %
When it is clear from the context, we do not write down the weakening action at all to
avoid clutter. 

Then we give the laws of the generic equivalence. %
Since the generic equivalence at layer $\metalevel$ subsumes that at layer $\proglevel$, we first give
the laws that hold for both layers, and then incrementally add those that only hold at
layer $\metalevel$.

\begin{law}[Subsumption]$ $
  \begin{itemize}
  \item If $\lttypgneeq i V{V'} l$, then $\lttypgeq i V {V'} l$. 
  \item If $\lttypgeq i T{T'} l$, then $\lttypeq i T{T'} l$. 
  \item If $\lttrmgneeq i \nu{\nu'} T l$, then $\lttrmgeq i \nu{\nu'} T l$.
  \item If $\lttrmgeq i t{t'} T l$, then $\lttyequiv i t{t'} T l$.
  \end{itemize}
\end{law}

As a lemma, subsumption propagates to contexts and local substitutions:
\begin{lemma}[Subsumption] $ $
  \begin{itemize}
  \item If $\ltctxgeq i \Gamma \Delta$, then $\lpequiv i \Gamma \Delta$. 
  \item If $\ltsubgeq i \delta {\delta'} \Delta$, then $\ltsubeq i \delta {\delta'} \Delta$.  
  \end{itemize}
\end{lemma}

Due to subsumption, we know that components in generic equivalence are well-formed or
well-typed:
\begin{lemma}[Presupposition] $ $
  \begin{itemize}
  \item If $\lttypgeq i T{T'} l$, then $\lttypwf i T l$ and $\lttypwf i{T'} l$.
  \item If $\lttrmgeq i t{t'} T l$, then $\lttyping i t T l$ and $\lttyping i{t'}T
    l$. 
  \end{itemize}
\end{lemma}
\begin{proof}
  By subsumption and presupposition. 
\end{proof}

\begin{law}[PER]
  All four relations are PERs. 
\end{law}

\begin{law}[Type Conversion] $ $
  \begin{itemize}
  \item If $\lttrmgneeq i \nu {\nu'} T l$ and $\lttypeq i T{T'} l$, then $\lttrmgneeq i \nu {\nu'}{T'} l$.
  \item If $\lttrmgeq i t {t'} T l$ and $\lttypeq i T{T'} l$, then $\lttrmgeq i t
    {t'}{T'} l$.
  \end{itemize}
\end{law}

\begin{law}[Context Equivalence] $ $
  \begin{itemize}
  \item If $\lttypgneeq i V {V'} l$, $\gequiv\Phi\Psi$ and
    $\lpequiv[\Phi]{i}\Delta\Gamma$, then $\lttypgneeq[\Phi][\Delta] i V {V'} l$.
  \item If $\lttypgeq i T{T'} l$, $\gequiv\Phi\Psi$ and
    $\lpequiv[\Phi]{i}\Delta\Gamma$, then $\lttypgeq[\Phi][\Delta] i T{T'} l$.
  \item If $\lttrmgneeq i \nu {\nu'} T l$, $\gequiv\Phi\Psi$ and
    $\lpequiv[\Phi]{i}\Delta\Gamma$, then $\lttrmgneeq[\Phi][\Delta] i \nu {\nu'}{T} l$.
  \item If $\lttrmgeq i t {t'} T l$, $\gequiv\Phi\Psi$ and
    $\lpequiv[\Phi]{i}\Delta\Gamma$, then $\lttrmgeq[\Phi][\Delta] i t {t'}{T} l$.
  \end{itemize}
\end{law}

\begin{law}[Weakening] $ $
  \begin{itemize}
  \item If $\lttypgneeq i V {V'} l$ and
    $\psi :: L \sep \Psi; \Gamma \To_i L' \sep \Phi;\Delta$, then
    $\lttypgneeq[\Phi][\Delta][L'] i{V}{V'}{l}$.
  \item If $\lttypgeq i T{T'} l$ and
    $\psi :: L \sep \Psi; \Gamma \To_i L' \sep \Phi;\Delta$, then
    $\lttypgeq[\Phi][\Delta][L'] i{T}{T'}{l}$.
  \item If $\lttrmgneeq i \nu {\nu'} T l$ and
    $\psi :: L \sep \Psi; \Gamma \To_i L' \sep \Phi;\Delta$, then
    $\lttrmgneeq[\Phi][\Delta][L'] i{\nu}{\nu'}{T}{l}$.
  \item If $\lttrmgeq i t {t'} T l$ and
    $\psi :: L \sep \Psi; \Gamma \To_i L' \sep \Phi;\Delta$, then
    $\lttrmgeq[\Phi][\Delta][L'] i{t}{t'}{T}{l}$.
  \end{itemize}  
\end{law}

\begin{law}[Weak Head Closure] $ $
  \begin{itemize}
  \item If $\ttypreds i T W l$, $\ttypreds i{T'}{W'} l$ and $\lttypgeq i W{W'} l$, then $\lttypgeq i T{T'} l$. 
  \item If $\ttrmreds i t w T l$, $\ttrmreds i{t'}{w'}T l$ and $\lttrmgeq i w{w'} T l$, then $\lttrmgeq i t{t'} T l$.
  \end{itemize}
\end{law}

\begin{law}[Type Constructors]
  If $\lpjudge i \Gamma$, 
  \begin{itemize}
  \item if $\tyequiv[L] l{l'}\Level$, then $\lttypgeq i{\Ty l}{\Ty l'}{1 + l}$ and
    $\lttrmgeq i{\Ty l}{\Ty l'}{\Ty{1 + l}}{2 + l}$;
  \item then $\lttypgeq i{\Nat}{\Nat}0$ and $\lttrmgeq i{\Nat}{\Nat}{\Ty 0}{1}$;
  \item if $\lttypgeq i S{S'} l$ and
    $\lttypgeq[\Psi][\Gamma, x : S \at l] i T{T'}{l'}$, then \newline
    $\lttypgeq i{\PI{l}{l'}x S T}{\PI{l}{l'}x{S'}{T'}}{l \sqcup l'}$;
  \item if $\lttrmgeq i s{s'}{\Ty l}{1 + l}$ and $\lttrmgeq[\Psi][\Gamma, x : \Elt l s \at l] i
    t{t'}{\Ty{l'}}{1 + l'}$, then \newline
    $\lttrmgeq i{\PI{l}{l'}x s t}{\PI{l}{l'}x{s'}{t'}}{\Ty{l \sqcup
      l'}}{1 + {(l \sqcup l')}}$.
  \end{itemize}
\end{law}

\begin{law}[Neutral Types]  $ $
  \begin{itemize}
  \item If $\lpjudge i \Gamma$, $U : \DTyp[\Delta]{i'} l \in \Psi$, $i' \in \{\codelevel, \proglevel\}$,
    $i' \le i$ and
    $\ltsubgeq i \delta{\delta'} \Delta$, then
    $\lttypgneeq i{U^\delta}{U^{\delta'}}{l}$. 
  \item If $\tyequiv[L]{l}{l'}\Level$ and $\lttrmgneeq i \nu{\nu'}{\Ty l}{1 + l}$,
    then $\lttypgneeq i{\Elt l \nu}{\Elt{l'}{\nu'}}{l}$.
  \end{itemize}
\end{law}

\begin{law}[Congruence] $ $
  \begin{itemize}
  \item If $\lpjudge i \Gamma$, then $\lttrmgeq i \ze \ze \Nat 0$.
  \item If $\lttrmgeq i t{t'}\Nat 0$, then $\lttrmgeq i {\su t}{\su{t'}}\Nat 0$.
  \item If $\lttypwf i S l$, $\lttyping i t {\PI{l}{l'}x{S}{T}}{l \sqcup l'}$,
    $\lttyping i {t'}{\PI{l}{l'}x{S}{T}}{l \sqcup l'}$ and \newline
    $\lttrmgeq[\Psi][\Gamma, x : S \at{l}] i
    {\APP{t}{l}{l'}x{S}{T}{x}}{\APP{t'}{l}{l'}x{S}{T}{x}}{T}{l'}$, then \newline
    $\lttrmgeq i t{t'}{\PI{l}{l'}x{S}{T}}{l \sqcup l'}$.
  \end{itemize}
\end{law}

\begin{law}[Congruence for Neutrals] $ $
  \begin{itemize}
  \item If $\lpjudge i \Gamma$ and $x : T \at l \in \Gamma$, then $\lttrmgneeq i x x
    T l$. 
  \item If $\lpjudge i \Gamma$, $u : \DTrm[\Delta]{i'}T l \in \Psi$,
    $i' \in \{\varlevel, \codelevel\}$, $i' \le i$ and $\ltsubgeq i \delta{\delta'} \Delta$, then
    \newline $\lttrmgneeq i{u^\delta}{u^{\delta'}} {T[\delta]} l$.
  \item If $\tyequiv[L]l{l'}\Level$, $\lttypgeq[\Psi][\Gamma, x : \Nat \at 0] i M{M'} l$,
    $\lttrmgeq i {s_1}{s_3} {M[\ze/x]}l$, \newline
    $\lttrmgeq[\Psi][\Gamma, x : \Nat \at 0, y : M \at l] i{s_2}{s_4}{M[\su
      x/x]}l$ and
    $\lttrmgneeq i \nu{\nu'} \Nat 0$, then \newline
    $\lttrmgneeq i{\ELIMN l{x.M}{s_1}{x,y. s_2}\nu}{\ELIMN{l'}{x.M'}{s_3}{x,y. s_4}{\nu'}}{M[\nu/x]}{l}$.
  \item If $\tyequiv[L]{l_1}{l_3}\Level$, $\tyequiv[L]{l_2}{l_4}\Level$, $\lttypgeq i S{S'}{l_1}$,
    $\lttypgeq[\Psi][\Gamma, x : S \at{l_1}] i{T}{T'}{l_2}$, 
    $\lttrmgneeq i \nu{\nu'}{\PI {l_1}{l_2} x{S''}{T''}}{l_1 \sqcup l_2}$ and $\lttrmgeq i
    s{s'} S{l_1}$, then \newline
    $\lttrmgneeq i{\APP \nu {l_1}{l_2} x S T s}{\APP{\nu'}{l_3}{l_4} x
      {S'}{T'}{s'}}{T[s/x]}{l_2}$. 

  \end{itemize}
\end{law}

We derive that
\begin{lemma}[Reflexivity of Local Identity Substitutions]
  If $\lpequiv i \Gamma \Delta$, then
  $\ltsubgeq i{\id_{\Gamma}}{\id_{\Gamma}}{\Delta}$.
\end{lemma}

\begin{lemma}[Congruence of Global Variables] $ $
  \begin{itemize}
  \item If $\lpequiv i \Gamma \Delta$, $u : \DTrm[\Delta]{i'}T l \in \Psi$, $i' \in \{\varlevel, \codelevel\}$ and
    $i' \le i$, then $\lttrmgneeq i {u^{\id_{\Gamma}}}{u^{\id_{\Gamma}}} T l$.
  \item If $\lpequiv i \Gamma \Delta$, $U : \DTyp[\Delta]{i'} l \in \Psi$,
    $i' \in \{\codelevel, \proglevel\}$ and $i' \le i$, then $\lttypgneeq i{U^{\id_{\Gamma}}}{U^{\id_{\Gamma}}}{l}$.
  \end{itemize}
\end{lemma}

At this point, we give the laws that should hold for both layers. %
Next, we consider laws that only hold at layer $i = \metalevel$. %
We first give the laws for type constructors.
\begin{law}[Type Constructors] $ $
  \begin{itemize}
  \item If $\lpjudge \metalevel \Gamma$, $\lttypgeq[\Psi][\Gamma][L, \vect\ell] \metalevel{T}{T'}{l}$
    and $\tyequiv[L,\vect\ell]{l}{l'}{\Level}$, then
    $\lttypgeq \metalevel{\UPI \ell l T}{\UPI \ell{l'}{T'}}\omega$.
  \item If $\lpjudge \metalevel \Gamma$, $\lttypgeq[\Psi, g : \Ctx] \metalevel{T}{T'}{l}$ and
    $\tyequiv[L]{l}{l'}{\Level}$, then \newline
    $\lttypgeq \metalevel{\CPI g l T}{\CPI g{l'}{T'}}{l}$. 
  \item If $\lpjudge \metalevel \Gamma$, $\lttypgeq[\Psi, U : \DTyp[\Delta] \proglevel{l_1}]
    \metalevel{T}{T'}{l_2}$, $\ltctxgeq \proglevel{\Delta}{\Delta'}$, $\tyequiv[L]{l_1}{l_3}{\Level}$
    and $\tyequiv[L]{l_2}{l_4}{\Level}$, then 
    $\lttypgeq \metalevel{\TPI U [\Delta]{l_1}{l_2} T}{\TPI U[\Delta']{l_3}{l_4}{T'}}{l_2}$.
  \item If $\lpjudge \metalevel \Gamma$ and $\ltctxgeq \proglevel{\Delta}{\Delta'}$, then $\lttypgeq
    \metalevel{\CTyp[\Delta]l}{\CTyp[\Delta']l}{0}$.
  \item If $\lpjudge \metalevel \Gamma$, $\ltctxgeq \proglevel{\Delta}{\Delta'}$ and
    $\lttypgeq[\Psi][\Delta]\proglevel T{T'}l$, then \newline
    $\lttypgeq \metalevel{\CTrm[\Delta] Tl}{\CTrm[\Delta']{T'}l}{0}$.
  \end{itemize}
\end{law}

\begin{law}[Congruence] $ $
  \begin{itemize}
  \item If $\lttyping \metalevel t {\UPI \ell l T}{\omega}$, $\lttyping \metalevel {t'}{\UPI \ell l
      T}{\omega}$ and $\lttrmgeq[\Psi][\Gamma][L, \vect\ell] \metalevel{\UAPP
      t{\ell}}{\UAPP{t'}{\ell}}{T}{l}$, then
    $\tconvtrmnf \metalevel{t}{t'}{\UPI \ell l T}{\omega}$. 
  \item If $\lttyping \metalevel t{\CPI g l T}{l}$, $\lttyping \metalevel {t'}{\CPI g l T}{l}$ and \newline
    $\lttrmgeq[\Psi, g : \Ctx] \metalevel{\CAPP t{g}}{\CAPP{t'}{g}}{T}{l}$, then $\lttrmgeq
    \metalevel{t}{t'}{\CPI g l T}{l}$. 
  \item If $\lttyping \metalevel t {\TPI U [\Delta]{l}{l'} T}{l'}$,
    $\lttyping \metalevel {t'}{\TPI U [\Delta]{l}{l'} T}{l'}$, and \newline
    $\lttrmgeq[\Psi, U : \DTyp[\Delta] \proglevel{l}] \metalevel{\TAPP
      t{U^{\id_\Delta}}}{\TAPP{t'}{U^{\id_\Delta}}}{T}{l'}$, then \newline
    $\lttrmgeq \metalevel{w}{w'}{\TPI U [\Delta]{l}{l'} T}{l'}$.
  \item If $\lpjudge \metalevel \Gamma$ and $\lttypwf[\Psi][\Delta] \codelevel T l$, then $\lttrmgeq
    \metalevel{\boxit T}{\boxit{T}}{\CTyp[\Delta] l}{0}$.
  \item If $\lpjudge \metalevel \Gamma$ and $\lttyping[\Psi][\Delta] \codelevel t T l$, then
    $\lttrmgeq \metalevel{\boxit t}{\boxit{t}}{\CTrm[\Delta] T l}{0}$. 
  \end{itemize}
\end{law}

\begin{law}[Congruence for Neutrals] $ $
  \begin{itemize}
  \item If $\lttrmgneeq \metalevel{\nu}{\nu'}{\UPI \ell l T}{\omega}$, $|\vect\ell| = |\vect l| = |\vect
    l'| > \codelevel$ and
    $\forall ~ \codelevel \le n < |\vect l| ~.~ \tyequiv[L]{\vect l(n)}{\vect l'(n)}\Level$,
    then $\lttrmgneeq \metalevel{\UAPP \nu l}{\nu'~\$~\vect l'}{T[\vect l/\vect \ell]}{l[\vect
      l/\vect \ell]}$. 
  \item If $\lttrmgneeq \metalevel{\nu}{\nu'}{\CPI g l T}{l}$ and $\ltctxgeq \proglevel
    \Delta{\Delta'}$, then $\lttrmgneeq \metalevel{\CAPP \nu
      \Delta}{\CAPP{\nu'}{\Delta'}}{T[\Delta/g]}{l}$.
  \item If $\lttrmgneeq \metalevel{\nu}{\nu'}{\TPI U[\Delta] l{l'}{T''}}{l'}$ and
    $\lttypgeq[\Psi][\Delta] \proglevel{T}{T'}l$, then \newline
    $\lttrmgneeq \metalevel{\TAPP \nu{T}}{\TAPP{\nu'}{T'}}{T''[T/U]}{l'}$. 
  \item If $\lpjudge \metalevel \Gamma$, $\tyequiv[L]{l_1}{l_3}\Level$, 
    $\tyequiv[L]{l_2}{l_4}\Level$, $\ltctxgeq \proglevel \Delta{\Delta'}$, \newline
    $\lttrmgneeq \metalevel{\nu}{\nu'}{\CTyp[\Delta]{l_2}}{0}$, 
    $\lttypgeq[\Psi][\Gamma,x_T : \CTyp[\Delta]{l_2} \at{0}]\metalevel{M}{M'}{l_1}$
    and \newline
    $ \lttrmgeq[\Psi, U : \DTyp[\Delta]\codelevel{l_2}]\metalevel{t_1}{t_2}{M[\boxit{U^\id}/x_T]}{l_1}$, then \newline
    $\lttrmgneeq \metalevel{\LETBTYP{l_1}{l_2}
      \Delta{x_T.M}{U}{t_1}\nu}{\LETBTYP{l_3}{l_4}{\Delta'}{x_T.M'}{U}{t_2}{\nu'}}{M[t/x_T]}{l_1}$. 
  \item If $\lpjudge \metalevel \Gamma$, $\tyequiv[L]{l_1}{l_3}\Level$, 
    $\tyequiv[L]{l_2}{l_4}\Level$, $\ltctxgeq \proglevel \Delta{\Delta'}$,
    $\lttypgeq \proglevel T{T'}{l_2}$,
    $\lttrmgneeq \metalevel{\nu}{\nu'}{\CTrm[\Delta]T{l_2}}{0}$, 
    $\lttypgeq[\Psi][\Gamma,x_T : \CTrm[\Delta]T{l_2} \at{0}]\metalevel{M}{M'}{l_1}$
    and \newline
    $ \lttrmgeq[\Psi, u : \DTrm[\Delta]\codelevel T{l_2}]\metalevel{t_1}{t_2}{M[\boxit{u^\id}/x_t]}{l_1}$, then \newline
    $\lttrmgneeq \metalevel{\LETBTRM{l_1}{l_2}
      \Delta{T}{x_t.M}{u}{t_1}\nu}{\LETBTRM{l_3}{l_4}{\Delta'}{T'}{x_T.M'}{u}{t_2}{\nu'}}{M[t/x_t]}{l_1}$. 
  \end{itemize}
\end{law}

Then we consider the law for neutral forms for recursive principles for code. %
The law follows a similar line to the equivalence judgments and the convertibility
checking: the evaluation is blocked when the scrutinee is neutral or is a $\tbox$'ed
global variable. 
\begin{law}[Neutral Recursion on Code] $ $
  \begin{itemize}
  \item If all motives and branches are related by corresponding generic equivalence,
    moreover, $\tyequiv[L]{l}{l'}\Level$, $\ltctxgeq \proglevel{\Delta}{\Delta'}$ and
    $\lttrmgneeq \metalevel{\nu}{\nu'}{\CTyp[\Delta]{l}}{0}$, then \newline
    $\lttrmgneeq \metalevel{\ELIMTYP{l_1}{l_2}Mbl\Delta
      \nu}{\ELIMTYP{l_3}{l_4}{M'}{b'}{l'}{\Delta'}{\nu'}}{M_\Typ[l/\ell,\Delta/g,\nu/x_T]}{l_1}$. 
  \item If all motives and branches are related by corresponding generic equivalence,
    moreover, $\tyequiv[L]{l}{l'}\Level$, $\ltctxgeq \proglevel{\Delta}{\Delta'}$,
    $U : \DTyp[\Delta]{\codelevel} l \in \Psi$ and $\ltsubst \codelevel \delta \Delta$, then \newline
    $\lttrmgneeq \metalevel{\ELIMTYP{l_1}{l_2}Mbl\Delta{(\boxit{U^\delta})}}{\ELIMTYP{l_3}{l_4}{M'}{b'}{l'}{\Delta'}{(\boxit{U^\delta})}}{M_\Typ[l/\ell,\Delta/g,\boxit{U^\delta}/x_T]}{l_1}$.
  \item If all motives and branches are related by corresponding generic equivalence,
    moreover, $\tyequiv[L]{l}{l'}\Level$, $\ltctxgeq \proglevel{\Delta}{\Delta'}$,
    $\tconvtyp \proglevel T{T'}{l}$ and
    $\tconvtrmnee \metalevel{\nu}{\nu'}{\CTrm[\Delta]T{l}}{0}$, then \newline
    $\lttrmgneeq \metalevel{\ELIMTRM{l_1}{l_2}Mbl\Delta T \nu}{\ELIMTYP{l_3}{l_4}{M'}{b'}{l'}{\Delta'}{T'}{\nu'}}{M_\Trm[l/\ell,\Delta/g,T/U_T,\nu/x_t]}{l_2}$.
  \item If all motives and branches are related by corresponding generic equivalence,
    moreover, $\tyequiv[L]{l}{l'}\Level$, $\ltctxgeq \proglevel{\Delta}{\Delta'}$,
    $\tconvtyp \proglevel T{T'}{l}$, $u : \DTrm[\Delta]{i'}T l \in \Psi$,
    $ i' \in \{\varlevel, \codelevel\}$ and $\ltsubst \codelevel \delta \Delta$, then \newline
    $\lttrmgneeq \metalevel{\ELIMTRM{l_1}{l_2}Mbl\Delta T
      {(\boxit{u^\delta})}}{\ELIMTYP{l_3}{l_4}{M'}{b'}{l'}{\Delta'}{T'}{(\boxit{u^\delta})}}{M_\Trm[l/\ell,\Delta/g,T/U_T,\boxit{u^\delta}/x_t]}{l_2}$.
  \end{itemize}
\end{law}
We conclude all the laws here fore the generic equivalence.

This generic equivalence will be instantiated twice times: the syntactic equivalence
judgments, and the convertibility checking judgments. %
\delamlang is way more complex than \citet{abel_decidability_2017}'s work because
\delamlang has computations at two layers, $\proglevel$ and $\metalevel$. %
Therefore, we must derive the necessary properties from the fundamental theorems to be
proved shortly at each layer. %

\subsection{Kripke Logical Relations for MLTT}\labeledit{sec:dt:klogrel}

The Kripke logical relations are parameterized by the generic equivalence. %
It is additionally indexed by another layering index $j \in \{\proglevel, \metalevel\}$, which
quantifies the types described by the relations. %
When $j = \proglevel$, we consider types from MLTT. %
When $j = \metalevel$, we consider all possible types. %
The reason for this distinction is to handle the lifting property from layer $\codelevel$
(which has the terms as $\proglevel$) or $\proglevel$ to
$\metalevel$, where terms from MLTT are brought to \delamlang. %
On the semantic side, we need to make sure that terms from MLTT can interact with
``native'' terms in \delamlang coherent. %
We further restrict $j = \proglevel$ when $i = \proglevel$. 

The Kripke logical relations are defined by
\begin{enumerate}
\item recursion on $j$, which effectively means the logical relations are
  2-layered; also note when $i = \proglevel$, $j = \proglevel$ is determined;
\item a transfinite well-founded recursion on the universe levels, and
\item induction-recursion on related types and terms. 
\end{enumerate}
In particular, the recursion on $j$ is necessary, as the relations when $j = \metalevel$ depend
on the validity judgments of $j = \proglevel$. %
When we do a recursion on universe levels, we must mind the well-foundedness of
universe levels. %
As we have discussed in \Cref{sec:dt:ulevel}, we are sure that all universe
levels must find a finite number of steps to descend to $0$. %
The only problem is $\omega$, which is not finite. %
Thus we must include one large cardinal to handle this level, hence the transfinite
recursion. %
Luckily, we do not have to think about it most of the time as we cannot really use
$\omega$ to do anything special at all. %
Note that our relations do not exactly follow \citet{abel_decidability_2017} tightly,
where two relations are defined for types and terms respectively. %
In our case, we provide simpler inductive-recursive definitions, where only one
relation is defined for types and for terms respectively. %
This style is more akin to the PER models in untyped domains by
\citet{abel_normalization_2013,hu_jang_pientka_2023}, except that our logical
relations are Kripke. %
We follow a proof schema that combines that of \citet{abel_decidability_2017} and that
of \citet{abel_normalization_2013,hu_jang_pientka_2023}. %
We define the following judgments:
\begin{itemize}
\item $\D :: \ldtypeq i j T{T'} l$ denotes that two types $T$ and $T'$ are related. %
  This relation is defined inductively. %
  We use $\D$ to mark give a name to the derivation as we will do recursion on it. %
\item $\ldtyequiv i j t{t'} \D$ denotes that two terms $t$ and $T'$ related by $\D$. %
  This relation is defined by a recursion on $\D$. 
\item $\E :: \ldctxeq i j \Gamma \Delta$ denotes that two contexts are related. %
  It is a generalization of $\D$. 
\item $\ldsubeq i j \delta{\delta'} \E$ denotes that two local substitutions
  $\delta$ and $\delta'$ are related. %
  It is a generalization of $\ldtyequiv i j t{t'} \D$ by doing recursion on $\E$. 
\end{itemize}
% The index $i$ is fixed for each instance of generic equivalence.
For convenience, we define the following:
\begin{align*}
  \ldtypwf i j T l &:= \ldtypeq i j T T l \\
  \ldtyequivt i j t{t'} T l & := \text{for some $T'$, } \D :: \ldtypeq i j T{T'}
l \tand \ldtyequiv i j t{t'}{\D} \\
  \ldtyping i j t T l & := \ldtyequivt i j t{t} T l \\
  \ldctxwf i j \Gamma &:=  \ldctxeq i j \Gamma \Gamma \\
  \ldsubeq i j \delta{\delta'} \Delta &:= \text{for some $\Delta'$, }\E :: \ldctxeq i
                                        j \Delta{\Delta'} \tand \ldsubeq i j \delta{\delta'} \E \\
  \ldsubst i j \delta \Delta &:= \ldsubeq i j \delta{\delta} \Delta
\end{align*}

Now we proceed to define the relations. %
We begin with the natural numbers. %
\begin{mathpar}
  \D :: \inferrule
  {\ttypreds i T{\Nat}0 \\ \ttypreds i {T'}\Nat0}
  {\ldtypeq i j T{T'}0}
\end{mathpar}
Then $\ldtyequiv i j t{t'} \D$ is defined by $\ldtyeqnat i t{t'}$, which we
define as follows:
\begin{mathpar}
  \inferrule
  {\ttrmreds i t w \Nat 0 \\ \ttrmreds i{t'}{w'}\Nat 0 \\ \lttrmgeq i w{w'}\Nat 0
    \\ \ldtynfeqnat i w{w'}}
  {\ldtyeqnat i t{t'}}

  \inferrule
  { }
  {\ldtynfeqnat i \ze \ze}

  \inferrule
  {\ldtyeqnat i t{t'}}
  {\ldtynfeqnat i{\su t}{\su{t'}}}

  \inferrule
  {\lttrmgneeq i \nu{\nu'}\Nat 0}
  {\ldtynfeqnat i{\nu}{\nu'}}  
\end{mathpar}

Then we consider universes.
\begin{mathpar}
  \D :: \inferrule
  {\ttypreds i T{\Ty{l_1}}{1 + l_1} \\ \ttypreds i {T'}{\Ty{l_2}}{1 + l_2} \\\\
  \tyequiv[L]{l_1}{l}\Level \\ \tyequiv[L]{l_2}{l}\Level}
  {\ldtypeq i j{T}{T'}{1 + l}}
\end{mathpar}
Then $\ldtyequiv i j t{t'} \D$ is defined by
\begin{itemize}
\item $\ttrmreds i t w{\Ty l}{1 + l}$,
\item $\ttrmreds i {t'}{w'}{\Ty l}{1 + l}$,
\item $\lttrmgeq i w{w'}{\Ty l}{1 + l}$, which means that $t$ and $t'$ are equivalent
  types at level $l$,
\item $\ldtypeq i j{\Elt l w}{\Elt l{w'}}{l}$, which means that the corresponding
  types of $w$ and $w'$ are related. 
\end{itemize}
The last condition requires the well-founded recursion on the universe levels in order
to refer back to the relation for types. %
Notice that the universe level decreases by one so this definition is valid. %

Then we define the relation for $\Pi$ types.
\begin{mathpar}
  \D :: \inferrule
  {\ttypreds i T{\PI l{l'} x{S_1}{T_1}}{l \sqcup l'} \\ \ttypreds i {T'}{\PI l{l'}
      x{S_2}{T_2}}{l \sqcup l'} \\
    \lttypwf i{S_1} l \\ \lttypwf i{S_2} l \\ \lttypwf[\Psi][\Gamma, x : S_1 \at l]
    i{T_1}{l'} \\
    \lttypwf[\Psi][\Gamma, x : S_2 \at l] i{T_2}{l'} \\
    \lttypgeq i{\PI l{l'} x{S_1}{T_1}}{\PI l{l'} x{S_2}{T_2}}{l \sqcup l'} \\
    \E :: (\forall~ \psi :: L' \sep \Phi;\Delta \To_i L \sep \Psi; \Gamma ~.~
    \ldtypeq[\Phi][\Delta][L'] i j{S_1}{S_2}{l}) \\
    \F :: (\forall~ \psi :: L' \sep \Phi;\Delta \To_i L \sep \Psi; \Gamma \tand \ldtyequiv[\Phi][\Delta][L']
    i j {s}{s'}{\E(\psi)} ~.~
    \ldtypeq[\Phi][\Delta][L'] i j{T_1[s/x]}{T_2[s'/x]}{l'})}
  {\ldtypeq i j{T}{T'}{l \sqcup l'}}
\end{mathpar}
This case is particularly complex. %
Let us digest the premises one by one. %
First, we require that both types $T$ and $T'$ to reduce to some $\Pi$ types. %
The typing judgments require the components of the $\Pi$ types are well-formed. %
Further, the $\Pi$ types themselves are equivalent. %
Then $\E$ requires that $S_1$ and $S_2$ are related under any weakening. %
This makes the relation Kripke. %
$\F$ is similar but require $T_1$ and $T_2$ remain related given any two related terms
$s$ and $s'$. %
Furthermore, in reality, we should put down sufficient premises for equivalent
universes to allow syntactically different universes to appear in both $T$ and $T'$,
but due to the size of this rule, we only use $l$ and $l'$ here. %
We apply the same principle in other rules in the rest of this technical report.

Then $\ldtyequiv i j t{t'} \D$ is defined by
\begin{itemize}
\item $\ttrmreds i t w{\PI l{l'} x{S_1}{T_1}}{l \sqcup l'}$, 
\item $\ttrmreds i {t'}{w'}{\PI l{l'} x{S_2}{T_2}}{l \sqcup l'}$,
\item $\lttrmgeq i w{w'}{\PI l{l'} x{S_1}{T_1}}{l \sqcup l'}$,
\item $\lttrmgeq i w{w'}{\PI l{l'} x{S_2}{T_2}}{l \sqcup l'}$, which is duplicated to
  make symmetry a simpler property,
\item
  for all $\forall\ \psi :: L' \sep \Phi;\Delta \To_i L \sep \Psi; \Gamma$ and $\Ac ::
  \ldtyequiv[\Phi][\Delta][L'] i j {s}{s'}{\E(\psi)}$, moreover, we add
  equivalence assumptions to remove the effect of type annotations,
  $\lttypgeq[\Phi][\Delta][L'] i{S_1}{S_1'}l$, $\lttypgeq[\Phi][\Delta][L'] i{S_2}{S_2'}l$, \break
  $\lttypgeq[\Phi][\Delta, x : S_1 \at l][L'] i{T_1}{T_1'}{l'}$,
  $\lttypgeq[\Phi][\Delta, x : S_2 \at l][L'] i{T_2}{T_2'}{l'}$, then \break
  $\ldtyequiv[\Phi][\Delta][L'] i j{\APP w l{l'} x{S_1'}{T_1'} s}{\APP{w'}l{l'}
    x{S_2'}{T_2'}{s'}}{\F(\psi, \Ac)}$.
\end{itemize}

The next case is neutral types.
\begin{mathpar}
  \D :: \inferrule
  {\ttypreds i T{V}l \\ \ttypreds i{T'}{V'}l \\
    \lttypgneeq i V{V'} l}
  {\ldtypeq i j{T}{T'}{l}}
\end{mathpar}
Then $\ldtyequiv i j t{t'} \D$ is defined by
\begin{itemize}
\item $\ttrmreds i t \nu{V}{l}$,
\item $\ttrmreds i {t'}{\nu'}{V'}{l}$,
\item $\lttrmgneeq i \nu{\nu'}{V}{l}$,
\item $\lttrmgneeq i \nu{\nu'}{V'}{l}$.
\end{itemize}
Again, we duplicate $\nu \sim \nu'$ to make symmetry easy.

The last case possible for both layers is equivalence of universe levels. %
This case is introduced for bureaucratic purposes and often ignored.
\begin{mathpar}
  \D :: \inferrule
  {\E :: \ldtypeq i j{T}{T'}{l'} \\ \tyequiv[L]l{l'}\Level}
  {\ldtypeq i j{T}{T'}{l}}
\end{mathpar}
Then $\ldtyequiv i j t{t'} \D$ is defined by $\ldtyequiv i j t{t'} \E$. 

Now we have given all logical relations for types and terms that are available at both layers. %
In order to give the logical relations for layer $\metalevel$, we need to first give the
logical relations for local contexts and local substitutions, as the semantics of
contextual types depend on them. %
The complexity of our logical relations primarily comes from the fact that we must be
able to handle computation at two layers ($\proglevel$ and $\metalevel$) and different lifting
behaviors. %

We proceed by defining related local contexts inductively and then the corresponding
recursive case for local substitutions. %
\begin{mathpar}
  \E :: \inferrule
  {\judge[L] \Psi}
  {\ldctxeq i j \cdot\cdot}
\end{mathpar}
Then $\ldsubeq i j \delta{\delta'} \E$ is defined by $\lpjudge i \Gamma$ and then checking $\Gamma$:
\begin{itemize}
\item if $\Gamma$ ends with $\cdot$, then $\delta = \delta' = \cdot^{|\Gamma|}$;
\item if $\Gamma$ ends with $g$, then $g : \Ctx \in \Psi$ and $\delta = \delta' = \cdot^{|\Gamma|}_g$;
\end{itemize}

\begin{mathpar}
  \E :: \inferrule
  {\judge[L] \Psi \\ g : \Ctx \in \Psi}
  {\ldctxeq i j {g}{g}}
\end{mathpar}
Then $\ldsubeq i j \delta{\delta'} \E$ is defined as $\lpjudge i \Gamma$ and
\begin{itemize}
\item $\Gamma$ also ends with $g$,
\item $\delta = \delta' = \wk^{|\Gamma|}_g$.
\end{itemize}

\begin{mathpar}
  \E :: \inferrule
  {\F :: \forall~\alpha :: L' \sep \Phi \To L \sep \Psi ~.~ \ldctxeq[\Phi][L'] i j
    \Delta{\Delta'} \\ \D :: \ldStypeq \F i j T{T'} l \\ \tyequiv[L]l{l'}\Level}
  {\ldctxeq i j{\Delta, x : T \at l}{\Delta', x : T' \at{l'}}}
\end{mathpar}
Then $\ldsubeq i j \delta{\delta'} \E$ is defined as
\begin{itemize}
\item $\delta = \delta_1, t/x$,
\item $\delta' = \delta_1', t'/x$,
\item $\Cc :: \forall~\alpha :: L' \sep \Phi \To L \sep \Psi ~.~ \ldsubeq[\Phi][\Gamma][L'] i j {\delta_1}{\delta_1'}{\F(\alpha)}$,
\item $\forall~\alpha :: L' \sep \Phi \To L \sep \Psi ~.~ \ldtyequiv[\Phi][\Gamma][L'] i j t{t'}{\D(\alpha,
    \Cc(\alpha))}$.
\end{itemize}
where we let $\D :: \ldStypeq \F i j T{T'} l$ to be
\[
  \forall~\alpha :: L' \sep \Phi \To L \sep \Psi \tand \ldsubeq[\Phi][\Gamma][L'] i j \delta{\delta'}{\F(\alpha)} ~.~ \ldtypeq[\Phi][\Gamma][L'] i j {T[\delta]}{T'[\delta']}{l}
\]
The judgment $\ldStypeq \F i j T{T'} l$ requires the stability under local substitutions
of the relation between $T$ and $T'$. %
We apply the same principle to related terms and derive the judgment
$\ldStyequiv \F i j t{t'}{\D}$, which is given by
\[
  \forall~\alpha :: L' \sep \Phi \To L \sep \Psi \tand \Cc :: \ldsubeq[\Phi][\Gamma][L'] i j
  \delta{\delta'}{\F(\alpha)} ~.~ \ldtyequiv[\Phi][\Gamma][L'] i j {t[\delta]}{t'[\delta']}{\D(\alpha,
    \Cc)}
\]

\subsection{Properties for Logical Relations When $j = \proglevel$}\labeledit{sec:dt:logrel-prop}


In this section, we pause our progress of defining the logical relations when $j =
\metalevel$. %
Following previous lines of work, when we give the Kripke logical relations to
contextual types, we will have to refer to the validity judgments of types and terms
in the logical relations. %
Therefore, in this section, we first work out a list of properties of the logical
relations when $j = \proglevel$, and then in the next section, we define semantic judgments for
global contexts and global substitutions, and then validity judgments. %
Once we have the validity judgments when $j = \proglevel$, we can then finish writing down the
logical relations for contextual types. %
Without further ado, let us start proving some lemmas. %
The list of lemmas mainly follows \citet{abel_normalization_2013,hu_jang_pientka_2023}
though we also take \citet{abel_decidability_2017} into consideration. 

\begin{lemma}[Weakening] $ $
  \begin{itemize}
  \item If $\D :: \ldtypeq i \proglevel{T}{T'}{l}$ and
    $\psi :: L' \sep \Phi;\Delta \To_i L \sep \Psi; \Gamma$, then
    $\E :: \ldtypeq[\Phi][\Delta][L'] i \proglevel{T}{T'}{l}$.
  \item If $\ldtyequiv i \proglevel{t}{t'}\D$ and
    $\psi :: L' \sep \Phi;\Delta \To_i L \sep \Psi; \Gamma$, then
    $\ldtyequiv[\Phi][\Delta][L'] i \proglevel{t'}{t}\E$.
  \end{itemize}
\end{lemma}
\begin{proof}
  Induction on $\D$. %
  Note that typing judgments are invariant under weakenings. %
  Also use the weakening law for generic equivalence. %
\end{proof}

\begin{lemma}[Escape] $ $
  \begin{itemize}
  \item If $\D :: \ldtypeq i \proglevel{T}{T'}{l}$, then $\lttypgeq i {T}{T'} l$.
  \item If $\ldtyequiv i \proglevel{t}{t'}\D$, then $\lttrmgeq i t{t'} T l$ and $\lttrmgeq i t{t'}{T'} l$.
  \end{itemize}  
\end{lemma}
\begin{proof}
  Induction on $\D$. %
  Use the subsumption law of the generic equivalence. %
\end{proof}

\begin{lemma}[Reflexivity of Neutral]
  If $\D :: \ldtypeq i \proglevel{T}{T'}l$, $\lttrmgneeq i \nu{\nu'} T l$ and $\lttrmgneeq i \nu{\nu'}{T'} l$, then $\ldtyequiv i \proglevel{\nu}{\nu'}\D$.
\end{lemma}
\begin{proof}
  Induction on $\D$. %
  Use the conversion and subsumption laws of generic equivalence and then the
  soundness of multi-step reductions and the escape lemma to obtain the equivalence
  between $\nu$ and $\nu'$ as normal forms. %
  In the function case, we use the congruence for neutral law of generic equivalence
  to show that the results of applying two neutral function are related.
\end{proof}

\begin{lemma}[Weak Head Expansion] $ $
  \begin{itemize}
  \item If $\D :: \ldtypeq i \proglevel{T}{T'}{l}$, $\ttypreds i{T_1} T l$ and $\ttypreds
    i{T_1'}{T'}l$, then \newline $\ldtypeq i \proglevel {T_1}{T_1'} l$.
  \item If $\ldtyequiv i \proglevel{t}{t'}\D$, $\ttrmreds i{t_1}{t}{T} l$ and $\ttrmreds
    i{t_1'}{t'}{T'}l$, then $\ldtyequiv i \proglevel {t_1}{t_1'} \D$.
  \end{itemize}  
\end{lemma}
\begin{proof}
  Induction on $\D$. Use transitivity of reductions. 
\end{proof}

\begin{lemma}[Symmetry] $ $
  \begin{itemize}
  \item If $\D :: \ldtypeq i \proglevel{T}{T'}{l}$, then $\E :: \ldtypeq i \proglevel{T'}{T}{l}$.
  \item If $\ldtyequiv i \proglevel{t}{t'}\D$, then $\ldtyequiv i \proglevel{t'}{t}\E$. 
  \end{itemize}
\end{lemma}
\begin{proof}
  We do induction on $\D$. %
  Since our definition is designed with symmetry in mind, symmetry is rather immediate.
\end{proof}

\begin{lemma}[Right Irrelevance]
  If $\D :: \ldtypeq i \proglevel{T}{T'}{l}$, $\E :: \ldtypeq i \proglevel{T}{T''}{l}$ and
  $\ldtyequiv i \proglevel{t}{t'}\D$, then $\ldtyequiv i \proglevel{t}{t'}\E$.
\end{lemma}
\begin{proof}
  We do induction on $\D$ and then invert $\E$. %
  We consider the function case. %
  In this case, we have premises
  \begin{align*}
    \D_1 &:: (\forall~ \psi :: L' \sep \Phi;\Delta \To_i L \sep \Psi; \Gamma ~.~
    \ldtypeq[\Phi][\Delta][L'] i \proglevel{S_1}{S_2}{l_1}) \\
    \D_2 &:: (\forall~ \psi :: L' \sep \Phi;\Delta \To_i L \sep \Psi; \Gamma ~.~ \forall~ \ldtyequiv[\Phi][\Delta][L']
    i \proglevel {s}{s'}{\D_1(\psi)} ~.~
    \ldtypeq[\Phi][\Delta][L'] i \proglevel{T_1[s/x]}{T_2[s'/x]}{l_2}) \\
    \E_1 &:: (\forall~ \psi :: L' \sep \Phi;\Delta \To_i L \sep \Psi; \Gamma ~.~
    \ldtypeq[\Phi][\Delta][L'] i \proglevel{S_1}{S_2'}{l_1}) \\
    \E_2 &:: (\forall~ \psi :: L' \sep \Phi;\Delta \To_i L \sep \Psi; \Gamma ~.~ \forall~ \ldtyequiv[\Phi][\Delta][L']
    i \proglevel {s}{s'}{\E_1(\psi)} ~.~
    \ldtypeq[\Phi][\Delta][L'] i \proglevel{T_1[s/x]}{T_2'[s'/x]}{l_2})
  \end{align*}
  from $\D$ and $\E$. %
  Here $\D_1$ and $\D_2$ are the premises of $\D$, and likewise for $\E$. %
  By determinacy, we know that
  \[
    T \reds \PI{l_1}{l_2} x{S_1}{T_1}
  \]
  must be unique.  %
  Moreover, we also know
  \begin{gather*}
    T' \reds \PI{l_1}{l_2} x{S_2}{T_2} \\
    T'' \reds \PI{l_1}{l_2} x{S_2'}{T_2'}
  \end{gather*}
  The most difficult part is to show that given
  \begin{itemize}
  \item $\psi :: L' \sep \Phi;\Delta \To_i L \sep \Psi; \Gamma$,    
  \item $\Ac :: \ldtyequiv[\Phi][\Delta][L'] i \proglevel {s}{s'}{\E_1(\psi)}$
  \item $\lttypgeq[\Phi][\Delta][L'] i{S_1}{S_3}l$,
  \item $\lttypgeq[\Phi][\Delta][L'] i{S_2}{S_4}l$, 
  \item $\lttypgeq[\Phi][\Delta, x : S_1 \at{l_1}][L'] i{T_1}{T_3}{l_2}$, and
  \item $\lttypgeq[\Phi][\Delta, x : S_2' \at{l_1}][L'] i{T_2'}{T_4}{l_2}$,
  \end{itemize}
  then
  \[
    \ldtyequiv[\Phi][\Delta][L'] i \proglevel{\APP w{l_1}{l_2} x{S_3}{T_3} s}{\APP{w'}{l_1}{l_2}
      x{S_4}{T_4}{s'}}{\E_2(\psi, \Ac)}
  \]

  From $\ldtyequiv i \proglevel{t}{t'}\D$, together with
  \begin{itemize}
  \item $\Bc :: \ldtyequiv[\Phi][\Delta][L'] i \proglevel {s}{s'}{\D_1(\psi)}$ by IH,
  \item $\lttypgeq[\Phi][\Delta, x : S_2 \at{l_1}][L'] i{T_2}{T_4}{l_2}$, where we
    know $\lttypeq i{S_2}{S_2'}{l_1}$ by escape, the subsumption law of generic
    equivalence, transitivity and local context equivalence of syntactic and generic
    equivalence. %
  \end{itemize}
  we have
  \[
    \ldtyequiv[\Phi][\Delta][L'] i \proglevel{\APP w{l_1}{l_2} x{S_3}{T_3} s}{\APP{w'}{l_1}{l_2}
      x{S_4}{T_4}{s'}}{\D_2(\psi, \Bc)}
  \]
  By another IH, we have the goal.
\end{proof}

\begin{lemma}[Left Irrelevance]
  If $\D :: \ldtypeq i \proglevel{T'}{T}{l}$, $\E :: \ldtypeq i \proglevel{T''}{T}{l}$ and
  $\ldtyequiv i \proglevel{t}{t'}\D$, then $\ldtyequiv i \proglevel{t}{t'}\E$.
\end{lemma}
\begin{proof}
  Immediate by symmetry and right irrelevance. 
\end{proof}

The left and right irrelevance lemmas are called the irrelevance lemma. %
It says that the exact relation between types is not important as long as their normal
forms are related. %

\begin{lemma}[Reflexivity and Transitivity] $ $
  \begin{itemize}
  \item If $\D_1 :: \ldtypeq i \proglevel{T_1}{T_2}{l}$ and $\D_2 :: \ldtypeq i
    \proglevel{T_2}{T_3}{l}$, then $\D_3 :: \ldtypeq i \proglevel{T_1}{T_3}{l}$.
  \item If $\E :: \ldtypeq i \proglevel{T_1}{T_1}{l}$, $\ldtyequiv i \proglevel{t_1}{t_2}{\D_1}$ and
    $\ldtyequiv i \proglevel{t_2}{t_3}{\D_2}$, then $\ldtyequiv i \proglevel{t_1}{t_3}{\D_3}$.
  \item $\F :: \ldtypeq i \proglevel{T_1}{T_1}{l}$.
  \item If $\ldtyequiv i \proglevel{t_1}{t_2}{\D_1}$, then $\ldtyequiv i \proglevel{t_1}{t_1}{\F}$.
  \end{itemize}
\end{lemma}
\begin{proof}
  We do induction on $\D_1$ and then invert $\D_2$. %
  The function case is the most complex one. %
  We have the following premises:
  \begin{align*}
    \Ac_1 &:: (\forall~ \psi :: L' \sep \Phi;\Delta \To_i L \sep \Psi; \Gamma ~.~
    \ldtypeq[\Phi][\Delta][L'] i \proglevel{S_1}{S_2}{l_1}) \\
    \Ac_2 &:: (\forall~ \psi :: L' \sep \Phi;\Delta \To_i L \sep \Psi; \Gamma ~.~ \forall~ \ldtyequiv[\Phi][\Delta][L']
    i \proglevel {s}{s'}{\Ac_1(\psi)} ~.~
    \ldtypeq[\Phi][\Delta][L'] i \proglevel{S_1'[s/x]}{S_2'[s'/x]}{l_2}) \\
    \Bc_1 &:: (\forall~ \psi :: L' \sep \Phi;\Delta \To_i L \sep \Psi; \Gamma ~.~
    \ldtypeq[\Phi][\Delta][L'] i \proglevel{S_2}{S_3}{l_1}) \\
    \Bc_2 &:: (\forall~ \psi :: L' \sep \Phi;\Delta \To_i L \sep \Psi; \Gamma ~.~ \forall~ \ldtyequiv[\Phi][\Delta][L']
    i \proglevel {s}{s'}{\E_1(\psi)} ~.~
    \ldtypeq[\Phi][\Delta][L'] i \proglevel{S_2'[s/x]}{S_3'[s'/x]}{l_2})
  \end{align*}
  from $\D$ and $\E$. %
  Here $\Ac_1$ and $\Ac_2$ are the premises of $\D_1$, and $\Bc_1$ and $\Bc_2$ are from $\D_2$. %
  By determinacy, we know that
  \begin{gather*}
    T_1 \reds \PI{l_1}{l_2} x{S_1}{S_1'} \\
    T_2 \reds \PI{l_1}{l_2} x{S_2}{S_2'} \\
    T_3 \reds \PI{l_1}{l_2} x{S_3}{S_3'}
  \end{gather*}

  Now, we should construct the transitivity for types and terms at the same time to
  understand how this proof is going to check out. %
  First we let 
  \[
    \Cc_1 :: (\forall~ \psi :: L' \sep \Phi;\Delta \To_i L \sep \Psi; \Gamma ~.~
    \ldtypeq[\Phi][\Delta][L'] i \proglevel{S_1}{S_3}{l_1})
  \]
  be the result of IH on $\Ac_1$ and $\Bc_1$. %
  Then our goal is to show that 
  
  The most difficult part is to show that given
  \begin{itemize}
  \item $\psi :: L' \sep \Phi;\Delta \To_i L \sep \Psi; \Gamma$,    
  \item $\F :: \ldtyequiv[\Phi][\Delta][L'] i \proglevel {s}{s'}{\Cc_1(\psi)}$
  \item $\lttypgeq[\Phi][\Delta][L'] i{S_1}{S_4}l$,
  \item $\lttypgeq[\Phi][\Delta][L'] i{S_3}{S_5}l$, 
  \item $\lttypgeq[\Phi][\Delta, x : S_1 \at{l_1}][L'] i{S_1'}{S_4'}{l_2}$, and
  \item $\lttypgeq[\Phi][\Delta, x : S_3 \at{l_1}][L'] i{S_3'}{S_5'}{l_2}$,
  \end{itemize}
  then
  \begin{itemize}
  \item
    \[
      \Cc_2 :: \ldtypeq[\Phi][\Delta][L'] i \proglevel{S_1'[s/x]}{S_3'[s'/x]}{l_2}
    \]
  \item 
    \[
      \ldtyequiv[\Phi][\Delta][L'] i \proglevel{\APP w_1{l_1}{l_2} x{S_4}{S_4'} s}{\APP{w_3}{l_1}{l_2}
        x{S_5}{S_5'}{s'}}{\Cc_2}
    \]
    where $t_k \reds w_k$ for $k \in \{1, 2, 3\}$. 
  \end{itemize}
  Our plan is the following:
  \begin{enumerate}
  \item We first relate $S_1'[s/x]$ and $S_2'[s/x]$, resp.
    $\APP {w_1}{l_1}{l_2} x{S_4}{S_4'} s$ and $\APP {w_2}{l_1}{l_2} x{S_4}{S_4'} s$, which
    requires
    \[
      \ldtyequiv[\Phi][\Delta][L'] i \proglevel {s}{s}{\Ac_1(\psi)}
    \]
    which is derived from reflexivity and irrelevance;
  \item and then relate $S_2'[s/x]$ and $S_3'[s/x]$, resp.
    $\APP {w_2}{l_1}{l_2} x{S_4}{S_4'} s$ and $\APP {w_2}{l_1}{l_2} x{S_5}{S_5'}{s'}$,
    which requires
    \[
      \ldtyequiv[\Phi][\Delta][L'] i \proglevel {s}{s'}{\Bc_1(\psi)}
    \]
    which is immediate from irrelevance.
  \end{enumerate}
  The final missing piece is reflexivity. %
  For types, it is just a result from transitivity and symmetry. %
  For terms, it is reflexivity and transitivity of related types and then
  irrelevance. %
  
  This concludes the function case. 
\end{proof}

Next, we work on the properties for related local contexts and local substitutions. %
We first consider the built property of weakening. %
Weakening can be seen as a case of the Yoneda lemma, where we only depend on
composition of weakenings.
\begin{lemma}[Weakening] $ $
  \begin{itemize}
  \item If $\D :: \ldctxeq i \proglevel \Delta{\Delta'}$ and  $\alpha :: L' \sep \Phi \To L \sep
    \Psi$, then $\E :: \ldctxeq[\Phi][L'] i \proglevel \Delta{\Delta'}$.
  \item If $\ldsubeq i \proglevel{\delta}{\delta'}\D$, $\alpha :: L' \sep \Phi \To L \sep \Psi$
    and $\tau :: L' \sep \Phi; \Gamma' \To_i \Gamma$, then
    $\ldsubeq[\Phi][\Gamma'][L'] i \proglevel{\delta}{\delta'}\E$.
  \end{itemize}
\end{lemma}
\begin{proof}
  Induction on $\D$. %
  In the two base cases, we know $\judge[L']\Phi$ by weakening. %
  In the step case, we simply store the given weakenings in a composition. %
\end{proof}
Notice that the second statement is a little bit strengthened. %
It is possible due to local weakening of related types and terms. %

We should first prove the reflexivity between identity local substitutions before
proving the the escape lemma.
\begin{lemma}[Reflexive Local Weakenings]
  If $\lpequiv i{\Delta,\Gamma}{\Delta',\Gamma}$ and $\D :: \ldctxeq i \proglevel \Delta{\Delta'}$, then \break
  $\ldsubeq[\Psi][\Delta,\Gamma] i \proglevel{\wk^{|\Gamma|}_{\Delta}}{\wk^{|\Gamma|}_{\Delta}}\D$. 
\end{lemma}
\begin{proof}
  We do induction on $\D$ and consider the step case
  \begin{mathpar}
    \D :: \inferrule
    {\E :: \forall \alpha :: L' \sep \Phi \To L \sep \Psi ~.~ \ldctxeq[\Phi][L'] i \proglevel {\Delta}{\Delta'} \\
      \F :: \ldStypeq \E i \proglevel T{T'} l \\ \tyequiv[L]l{l'}\Level}
    {\ldctxeq i \proglevel{\Delta, x : T \at l}{\Delta', x : T' \at{l'}}}
  \end{mathpar}
  \begin{align*}
    \Bc :: & \forall \alpha :: L' \sep \Phi \To L \sep \Psi ~.~ \ldsubeq[\Phi][\Delta, x : T \at l,\Gamma][L'] i \proglevel{\wk^{1 + |\Gamma|}_{\Delta}}{\wk^{1 + |\Gamma|}_{\Delta}}{\E(\alpha)}
      \byIH \\
    \alpha :: & L' \sep \Phi \To L \sep \Psi
                \tag{by assumption} \\
    & \ldtypeq[\Phi][\Delta, x : T \at l,\Gamma][L'] i \proglevel {T[\wk^{1 + |\Gamma|}_{\Delta}]}{T'[\wk^{1 + |\Gamma|}_{\Delta}]}{l}
      \tag{as $\F(\alpha, \Bc(\alpha))$} \\
    \Ac :: & \ldtypeq[\Phi][\Delta, x : T \at l,\Gamma][L'] i \proglevel {T}{T'}{l} \\
    & \lpjudge[\Phi][L'] i{\Gamma, x : T \at l,\Delta}
      \tag{by weakening and presupposition} \\
    & \lttrmgneeq[\Phi][\Delta, x : T \at l,\Gamma][L'] i x x T l
      \tag{by congruence for neutrals of generic equivalence} \\
    & \ldtyequiv[\Phi][\Delta, x : T \at l,\Gamma][L'] i \proglevel{x}{x}\Ac
      \tag{by reflexivity of neutral} \\
    & \ldsubeq[\Psi][\Gamma, x : T \at l,\Delta] i \proglevel{\wk^{1 + |\Gamma|}_{\Delta},x/x}{\wk^{1 + |\Gamma|}_{\Delta},x/x}\D
  \end{align*}
  Thus we conclude the goal. 
\end{proof}

\begin{corollary}[Reflexive Local Identity Substitutions]\labeledit{lem:dt:refl-lid}
  If $\lpequiv i \Gamma \Delta$ and $\D :: \ldctxeq i \proglevel \Gamma \Delta$, then \break
  $\ldsubeq i \proglevel{\id}{\id}\D$. 
\end{corollary}
\begin{proof}
  This is a specialization of the previous lemma. 
\end{proof}

Knowing all local identity substitutions are reflexively related, we can then prove
the escape lemma.
\begin{lemma}[Escape] $ $
  \begin{itemize}
  \item If $\D :: \ldctxeq i \proglevel{\Delta_1}{\Delta_2}$, then $\ltctxgeq i{\Delta_1}{\Delta_2}$. 
  \item If $\ldsubeq i \proglevel{\delta}{\delta'}\D$, then $\ltsubgeq i
    {\delta}{\delta'}{\Delta_1}$ and $\ltsubgeq i {\delta}{\delta'}{\Delta_2}$.
  \end{itemize}
\end{lemma}
\begin{proof}
  Induction on $\D$. %
  We consider the step case.
  \begin{mathpar}
    \D :: \inferrule
    {\E :: \forall \alpha :: L' \sep \Phi \To L \sep \Psi ~.~ \ldctxeq[\Phi][L'] i
      \proglevel{\Delta_1}{\Delta_2} \\ \F :: \ldStypeq \E i \proglevel T{T'} l \\ \tyequiv[L]l{l'}\Level}
    {\ldctxeq i \proglevel{\Delta_1, x : T \at l}{\Delta_2, x : T' \at{l'}}}
  \end{mathpar}
  \begin{align*}
    & \ltctxgeq i{\Delta_1}{\Delta_2}
      \byIH \\
    & \ldsubeq[\Psi][\Delta_1] i \proglevel{\id}{\id}{\E(\id)}
      \tag{by \Cref{lem:dt:refl-lid}} \\
    & \ldtypeq[\Psi][\Delta_1] i \proglevel {T}{T'}{l}
      \tag{as $\F(\id, \E(\id))$} \\
    & \lttypgeq[\Psi][\Delta_1] i {T}{T'} l
      \tag{by escape}
  \end{align*}
  Hence we conclude the first statement.
  
  In the second statement, we have $\delta = \delta_1, t/x$ and $\delta' = \delta_1',
  t'/x$, then
  \begin{align*}
    & \ldsubeq i \proglevel {\delta_1}{\delta_1'}{\E(\id)}
      \tag{by assumption} \\
    & \ltsubgeq i {\delta_1}{\delta_1'}{\Delta_1} \tand \ltsubgeq i {\delta_1}{\delta_1'}{\Delta_2}
      \byIH \\
    & \ldtyequiv i \proglevel t{t'}{\F(\id, \E(\id))}
      \tag{by assumption} \\
    & \lttrmgeq i t{t'}{T[\delta_1]}l
      \tag{by escape of related terms}
  \end{align*}
  Therefore we conclude the second statement as well.
\end{proof}

\begin{lemma}[Symmetry] $ $
  \begin{itemize}
  \item If $\D :: \ldctxeq i \proglevel{\Delta_1}{\Delta_2}$, then $\E :: \ldctxeq i \proglevel{\Delta_2}{\Delta_1}$.
  \item If $\ldsubeq i \proglevel{\delta}{\delta'}\D$, then $\ldsubeq i \proglevel{\delta'}{\delta}\E$.
  \end{itemize}  
\end{lemma}
\begin{proof}
  Induction on $\D$. %
  Use the symmetry of related types to obtain the goal.
\end{proof}

\begin{lemma}[Right Irrelevance]
  If $\D :: \ldctxeq i \proglevel{\Delta_1}{\Delta_2}$, $\E :: \ldctxeq i
  \proglevel{\Delta_1}{\Delta_3}$ and $\ldsubeq i \proglevel{\delta}{\delta'}\D$, then
  $\ldsubeq i \proglevel{\delta}{\delta'}\E$.
\end{lemma}
\begin{proof}
  Do induction on $\D$ and then invert $\E$. %
  Use the right irrelevance of related terms. %
\end{proof}

\begin{lemma}[Left Irrelevance]
  If $\D :: \ldctxeq i \proglevel{\Delta_1}{\Delta_2}$, $\E :: \ldctxeq i
  \proglevel{\Delta_3}{\Delta_1}$ and $\ldsubeq i \proglevel{\delta}{\delta'}\D$, then
  $\ldsubeq i \proglevel{\delta}{\delta'}\E$.
\end{lemma}
\begin{proof}
   A direct consequence of right irrelevance and symmetry. 
\end{proof}

\begin{lemma}[Reflexivity and Transitivity] $ $
  \begin{itemize}
  \item If $\D_1 :: \ldctxeq i \proglevel{\Delta_1}{\Delta_2}$ and
    $\D_2 :: \ldctxeq i \proglevel{\Delta_2}{\Delta_3}$, then
    $\D_3 :: \ldctxeq i \proglevel{\Delta_1}{\Delta_3}$.
  \item If $\ldsubeq i \proglevel{\delta_1}{\delta_2}{\D_1}$ and
    $\ldsubeq i \proglevel{\delta_2}{\delta_3}{\D_2}$, then
    $\ldsubeq i \proglevel{\delta_1}{\delta_3}{\D_3}$.
  \item $\E :: \ldctxeq i \proglevel{\Delta_1}{\Delta_1}$.
  \item If $\ldsubeq i \proglevel{\delta_1}{\delta_2}{\D_1}$, then $\ldsubeq i \proglevel{\delta_1}{\delta_1}{\E}$.
  \end{itemize}  
\end{lemma}
\begin{proof}
  Do induction on $\D_1$ and then invert $\D_2$. %
  The proof proceeds very similarly to the relations for types and terms. %
  We only consider the step case. %
  In this case, we have the following premises:
  \begin{itemize}
  \item $\F_1 :: \forall~\alpha :: L' \sep \Phi \To L \sep \Psi ~.~
    \ldctxeq[\Phi][L'] i \proglevel {\Delta_1'}{\Delta_2'}$,
  \item $\F_2 :: \forall~\alpha :: L' \sep \Phi \To L \sep \Psi ~.~
    \ldctxeq[\Phi][L'] i \proglevel {\Delta_2'}{\Delta_3'}$,
  \item $\Ac_1 :: \ldStypeq{\F_1} i \proglevel {T_1}{T_2} l$,
  \item $\Ac_2 :: \ldStypeq{\F_2} i \proglevel{T_2}{T_3} l$,
  \item If $\ldStypeq{\F_1} i \proglevel T{T'} l$ and $\ldStypeq{\F_2} i \proglevel{T'}{T''} l$, then
    $\ldStypeq {\F_3} i \proglevel T{T''} l$.
  \item
    $\forall~\alpha :: L' \sep \Phi \To L \sep \Psi ~.~ \ldtyequiv[\Phi][\Gamma][L']
    i \proglevel {t_1}{t_2}{\Ac_1(\alpha, \F_1(\alpha))}$,
  \item
    $\forall~\alpha :: L' \sep \Phi \To L \sep \Psi ~.~ \ldtyequiv[\Phi][\Gamma][L']
    i \proglevel {t_2}{t_3}{\Ac_2(\alpha, \F_2(\alpha))}$,
  \end{itemize}
  By IH, it is easy to show
  \[
    \F_3 :: \forall~\alpha :: L' \sep \Phi \To L \sep \Psi ~.~
    \ldctxeq[\Phi][L'] i \proglevel {\Delta_1'}{\Delta_3'}
  \]
  The difficult goals are
  \begin{itemize}
  \item $\Ac_3 :: \ldStypeq{\F_3} i \proglevel{T_1}{T_3} l$, and
  \item
    $\forall~\alpha :: L' \sep \Phi \To L \sep \Psi ~.~ \ldtyequiv[\Phi][\Gamma][L']
    i \proglevel {t_1}{t_3}{\Ac_3(\alpha, \F_3(\alpha))}$. 
  \end{itemize}
  We first assume $\alpha :: L' \sep \Phi \To L \sep \Psi$. %
  To prove the first statement, we further assume $\ldsubeq[\Phi][\Gamma][L'] i j
  \delta{\delta'}{\F_3(\alpha)}$. %
  Then we do a similar reasoning to the function case for transitivity of related
  types. %
  \begin{enumerate}
  \item We first relate $T_1[\delta]$ and $T_2[\delta]$ by using reflexivity.
  \item Then we related $T_2[\delta]$ and $T_3[\delta']$.
  \item Then we apply transitivity of related types. 
  \end{enumerate}

  To prove the second statement, we can simply use irrelevance so that transitivity
  can eventually apply on $\Ac_3(\alpha, \F_3(\alpha))$. 
\end{proof}

\begin{lemma}[Transitivity]
  Given
  \begin{itemize}
  \item 
    $\F_1 :: \forall~\alpha :: L' \sep \Phi \To L \sep \Psi ~.~ \ldctxeq[\Phi][L'] i \proglevel
    {\Delta_1}{\Delta_2}$,
  \item $\F_2 :: \forall~\alpha :: L' \sep \Phi \To L \sep \Psi ~.~ \ldctxeq[\Phi][L'] i \proglevel
    {\Delta_2}{\Delta_3}$, and
  \item $\F_3 :: \forall~\alpha :: L' \sep \Phi \To L \sep \Psi ~.~ \ldctxeq[\Phi][L'] i \proglevel
    {\Delta_1}{\Delta_3}$,
  \end{itemize}
  then
  \begin{itemize}
  \item if $\D_1 :e: \ldStypeq{\F_1} i \proglevel T{T'} l$ and $\D_2 :: \ldStypeq{\F_2} i
    \proglevel{T'}{T''} l$, then 
    $\D_3 :: \ldStypeq{\F_3} i \proglevel T{T'} l$;
  \item if $\ldStyequiv{\F_1} i \proglevel t{t'}{\D_1}$ and $\ldStyequiv{\F_2} i \proglevel
    {t'}{t''}{\D_3}$, then $\ldStyequiv{\F_3} i \proglevel t{t''}{\D_3}$.
  \end{itemize}
\end{lemma}
\begin{proof}
  The first statement is what we have proved in transitivity previously. %
  The second statement uses the transitivity of related terms. %
\end{proof}

\input{sec_validp}

\input{sec_fund}
